\answer{ex:08:1}{(a)~$\approx1.7\cdot10^6$. (b)~$g=\frac{0.9}{\sqrt{0.19}}\,\sigma_{\text{depois}}/\sigma_{\text{antes}}\approx16.5$: a resposta só depende da razão entre os ruídos.}
\answer{ex:08:2}{(a)~Autovalores $-(1\pm g\beta)/\tau$; a diferença decai em $\tau/(1-g\beta)$: $40\ms$ com $g\beta=0.5$ (diferença de entradas amplificada três vezes), $200\ms$ com $0.9$ ($19$ vezes). (b)~$0.905$ e $1.105$; entre as fronteiras, só a entrada forte vence.}
\answer{ex:08:3}{$P(k)=e^{gu_k/\sigma}/\sum_le^{gu_l/\sigma}$, ou seja,~\eqref{eq:softmax}; $g=10\ln9\approx22$.}
\answer{ex:08:4}{$\bar\Delta=1/3$, $g^*\approx3\ln(1/h)$: $21$, $14$, $9$ contra $16$--$23$, $11$ e $8$ no modelo.}
\answer{ex:08:5}{$g^*$ de $16$--$23$ com $h=0.001$ a $6$--$8$ com $h=0.2$; a estimativa $3\ln(1/h)$ vale dentro de um fator $1.5$ para $h\le0.05$; com $h\gtrsim\eta/10$, a regra $Q\leftarrow Q+\eta(R-Q)$ não consegue assimilar o que a busca obteve, e $g^*$ empaca perto de $8$.}
\answer{ex:08:6}{$T_p=\tau\ln9=44\ms$ com $g_{\text{lo}}=0$ (modelo, $44\ms$) e $70\ms$ com $g_{\text{lo}}=0.25$: uma rajada de $2$--$3\tau$ com ganho perto de zero.}
\answer{ex:08:7}{(a)~Melhor constante $0.925$ ($g=8$), oráculo $0.929$ ($16/8$): quase não há o que ganhar. (b)~Por exemplo, épocas calmas com recompensas raras ($p\in[0,0.3]$) e instáveis com $h=0.1$: $0.850$ contra $0.872$; o ajuste captura cerca de um quarto com $\eta=0.2$ e quase tudo com $\eta=0.05$.}
